\documentclass[12pt,a4paper]{article}
\usepackage[margin=2.5cm]{geometry}
\usepackage{setspace}
\usepackage{graphicx}
\usepackage{float}
\usepackage{booktabs}
\usepackage{longtable}
\usepackage{array}
\usepackage{hyperref}
\usepackage{enumitem}
\usepackage{titlesec}
\usepackage{xeCJK}
\setCJKmainfont{Noto Serif CJK TC}
\setmainfont{Times New Roman}
\hypersetup{colorlinks=true,linkcolor=black,urlcolor=blue,citecolor=black}
\setstretch{1.35}
\titleformat{\section}{\Large\bfseries}{\thesection}{1em}{}
\titleformat{\subsection}{\large\bfseries}{\thesubsection}{1em}{}

\title{基於大語言模型的智能教學助手設計與實現\\
\large -- 以澳門高中數學「平面向量」單元為例}
\author{作者：蕭如瀚 \\ 指導老師：黃達森}
\date{2025 年}

\begin{document}
\maketitle
\thispagestyle{empty}
\newpage

\section*{摘要}
本研究針對澳門教育暨青年發展局（DSEDJ）高中數學「平面向量」單元，設計並實現一套基於大型語言模型（Large Language Model, LLM）的智能教學助手（Intelligent Tutoring System, ITS）Web 原型。系統定位為課後鞏固輔具，整合診斷問卷、動態題庫練習、蘇格拉底式（Socratic）引導、三層降級檢索增強生成（Retrieval-Augmented Generation, RAG）與掌握度視覺化儀表板，形成「診斷--練習--回饋--追蹤--調整」個性化學習閉環。

技術實作採用 Next.js App Router、React、TypeScript、Supabase（PostgreSQL + pgvector extension）、DeepSeek Chat Completions、KaTeX 與 MathLive。本文融合既有兩版論文的優點：保留 v2 的敘事清晰與章節緊湊，並吸收 v3 的工程可重現性與證據鏈細節（完整 Prompt 對照、API 契約、sessionStorage 鍵位、AI trace 記錄）。核心工程設計包括：（1）雙階段命題管線（Generate \& Validate），以錯誤原因碼降低自動命題風險；（2）Hint 層級推斷與 \texttt{UI\_SIGNAL} 協議，提升 Socratic 引導與前端狀態一致性；（3）RAG 三層降級策略（同 topic 取樣 $\rightarrow$ 向量 RPC $\rightarrow$ ILIKE 後備），提高小題庫情境下的穩定性；（4）全鏈路 JSONL trace，支援後續量化評估與故障歸因。

研究結果顯示，本系統已完成端到端功能實作，並在原型驗證層面展現良好可用性與可觀測性。惟教育成效仍需透過較大樣本、較長週期的對照研究進一步驗證。

\noindent\textbf{關鍵詞：}大語言模型、智能教學系統、平面向量、蘇格拉底式引導、RAG、可觀測性

\newpage
\tableofcontents
\newpage

\section{緒論}
\subsection{研究背景與問題陳述}
向量是高中數學中幾何與代數高度耦合的核心單元。澳門 DSEDJ 高中數學大綱要求學生同時掌握向量運算、向量模長、數量積、共線與垂直判斷、向量幾何應用等能力。相較於單一技能型題目，向量題常要求學生在幾何直觀與代數推導間持續切換，認知負荷較高。

依據實際教學觀察，學生常見困難包含：第一，步驟跳躍，面對板書省略中間推導時難以自主補全；第二，公式與概念脫節，能背誦公式但難以解釋其幾何意義；第三，多步推理脆弱，任一步驟出錯即導致整題失分。現有多數通用 ITS 或 LLM 教學工具雖可提供即時回應，但未必對齊澳門課綱內容與教學語境。

本研究因此定位於「本地課綱對齊 + 課後鞏固 + 低成本可部署 + 可重現工程」四項目標，建立可落地的智能教學原型。

\subsection{研究目標與核心問題}
本研究以系統設計與工程驗證為主軸，聚焦四個問題：
\begin{enumerate}[leftmargin=2em]
  \item 如何把診斷結果與練習行為轉化為可追蹤、可解釋的個性化回饋？
  \item 如何在延遲與成本可控下，降低 AI 自動命題的質量錯誤率？
  \item 如何避免 Socratic 引導模板化，並識別學生已掌握內容？
  \item 如何在小題庫與有限基礎設施下，實作可降級的 RAG 檢索策略？
\end{enumerate}

\subsection{研究貢獻}
\begin{itemize}[leftmargin=2em]
  \item \textbf{教學設計貢獻：}建立單題（微觀）--批次（中觀）--全局（宏觀）三層回饋閉環。
  \item \textbf{工程設計貢獻：}提出雙階段命題、Hint 層級推斷覆寫、RAG 三層降級策略。
  \item \textbf{可觀測性貢獻：}以 JSONL 記錄 prompt/response/result/error，形成可分析數據基礎。
  \item \textbf{重現性貢獻：}提供 Prompt 對照、契約欄位與部署流程，降低二次實作門檻。
\end{itemize}

\section{文獻與理論基礎}
\subsection{傳統 ITS 的價值與限制}
傳統 ITS 在結構化計算教學具備成熟成果，優勢在規則可控與評分穩定。然而當任務需要語言推理、概念遷移與多步驟解釋時，純規則系統容易在回饋深度與語言彈性上受限。此差異成為 LLM 融入教學系統的重要動機。

\subsection{LLM 在數學教育的機會與風險}
LLM 能提供自然語言互動與個性化解釋，有助於提升學習動機與即時支持能力；但同時存在推理幻覺、數值錯誤、對話一致性不足與過度依賴風險。因此，教學場景中的 LLM 應搭配結構化輸出、驗題機制與降級策略，而非單純依賴模型自由生成。

\subsection{Socratic 引導與提示工程}
Socratic 導向強調以提問促進學生主動建構。其實務難點在於：模型若缺乏狀態感知，容易重複同質問句或忽略學生已答對訊號。本文採取「提示詞規則 + 程式邏輯覆寫」雙軌設計，提升引導的情境一致性。

\section{需求分析與教學設計}
\subsection{課程邊界}
本系統聚焦平面向量，使用 13 個 topic\_tag（如 \texttt{vec\_magnitude}、\texttt{dot\_product}、\texttt{vec\_projection} 等）描述知識維度。難度採 L1--L4 分層，避免超綱內容干擾基礎教學目標。

\subsection{功能需求與學習閉環}
系統主要由四個互補模組構成：\texttt{/diagnostic}（診斷問卷）、\texttt{/practice}（練習與引導）、\texttt{/mastery}（掌握度儀表板）、\texttt{/api/*}（後端推理與編排）。學生流程為：完成診斷 $\rightarrow$ 接收結構化報告 $\rightarrow$ 進入動態練習與對話 $\rightarrow$ 取得批次與總結回饋。

\begin{figure}[H]
  \centering
  \includegraphics[width=0.92\textwidth]{figures/arch-overview.png}
  \caption{系統整體架構圖（預留：請替換為網站或架構圖截圖）}
  \label{fig:arch-overview}
\end{figure}

\subsection{資料與介面契約}
為支持無帳號快速部署，系統將診斷、練習事件、批次報告、掌握報告等存於 \texttt{sessionStorage}。導師端回覆採用「可讀文字 + 機器信號」雙層輸出：回覆末尾附 \texttt{<!-- UI\_SIGNAL -->} JSON，供前端決策 \texttt{show\_hint\_button}、\texttt{next\_question} 等動作。

\section{系統架構與技術選型}
\subsection{三層架構}
系統採前後端一體化單倉架構：
\begin{enumerate}[leftmargin=2em]
  \item \textbf{表現層：}React 頁面與互動組件（KaTeX 顯示、MathLive 輸入）。
  \item \textbf{應用層：}Next.js Route Handlers（業務邏輯、LLM 調用、RAG）。
  \item \textbf{資料層：}Supabase 題庫與可選向量檢索（pgvector）。
\end{enumerate}

\subsection{技術選型原則}
選型重點為：可部署性、低成本、TypeScript 型別安全、與 OpenAI-compatible API 的互操作。此設計便於在教學現場快速複製，並降低維運負擔。

\section{核心演算法與實作}
\subsection{診斷評分與報告生成}
診斷問卷以 L1=1、L2=2、L3=3 的加權規則計算掌握度，並以護欄規則抑制偶然高分。後端流程包括：填空複核、Prompt A 報告生成、欄位級 merge/fallback、topic\_tag 校正，確保輸出可解析且可用。

\subsection{雙階段命題（Generate \& Validate）}
命題端先生成候選題，再由獨立驗題 prompt 依五類錯誤碼審核：
\texttt{ANSWER\_MISMATCH}、
\texttt{OPTIONS\_DUPLICATE}、
\texttt{STEM\_OPTION\_CONTRADICTION}、
\texttt{UNSOLVABLE\_OR\_NO\_VALID\_CHOICE}、
\texttt{FORMAT\_INVALID}。
若未達目標數量，系統在上限輪次內補生；超限後回退題庫題以保證可用性。

\begin{figure}[H]
  \centering
  \includegraphics[width=0.92\textwidth]{figures/generate-validate-flow.png}
  \caption{雙階段命題流程圖（預留：可用網站流程截圖或自繪圖）}
  \label{fig:gen-validate}
\end{figure}

\subsection{Socratic 導師與 Hint 層級推斷}
\texttt{inferHintIndexFromThoughts} 將學生思路與提示文本正規化後比對關鍵詞覆蓋率，推斷已涵蓋的提示層。若推斷層級高於模型輸出，伺服端會覆寫 \texttt{hint\_index\_used} 並觸發 \texttt{auto\_reveal\_hint}，避免「學生已懂但系統仍重複追問」。

\begin{figure}[H]
  \centering
  \includegraphics[width=0.82\textwidth]{figures/tutor-chat-ui.png}
  \caption{向量老師對話頁（預留：網站實際截圖）}
  \label{fig:tutor-chat}
\end{figure}

\subsection{RAG 三層降級}
RAG 檢索分三層：Layer 1 同 topic 取樣（不依賴 embedding）、Layer 2 向量 RPC、Layer 3 ILIKE 文字後備。當 embedding 或 RPC 不可用時，系統可自動降級，不會中斷主流程。

\subsection{掌握報告與正規化}
\texttt{/api/mastery/report} 將診斷、練習事件與批次報告組裝後交由模型生成 \texttt{MasteryReportV2}，伺服端再進行欄位正規化（難度、topic 分數、順序去重、敘事長度）並提供 fallback，確保儀表板穩定渲染。

\section{Prompt 工程與可靠性設計}
\subsection{Prompt 分工策略}
本研究採「低溫結構化、高溫對話化」原則：0.10--0.25 用於排序/驗題/JSON 任務，0.30--0.35 用於生成與敘事，0.45 用於導師對話與提示改寫。此分工可同時兼顧可解析性與語言自然度。

\subsection{全鏈路 fallback}
每個核心端點皆設降級策略：無金鑰時回本地規則、JSON 失敗時吞異常並回 fallback、驗題超限時回題庫補齊、RAG 失效時降至後備層。此設計保障服務連續性。

\subsection{AI trace 可觀測性}
啟用 \texttt{AI\_TRACE\_LOG=1} 後，系統將 prompt/response/error/result 以 JSONL 追加記錄至 \texttt{.debug/ai-traces/}。trace 內建敏感鍵遮罩與長內容截斷，可用於 prompt 比對、錯誤碼統計與調參分析。

\section{系統驗證與評估設計}
\subsection{功能驗證結果}
以工程驗證角度，系統已完成診斷、練習、對話、批次報告、掌握儀表板與降級機制的端到端串接。包含「無 embedding」、「RPC 未建立」、「模型回傳不合規 JSON」等失效情境，皆可透過 fallback 維持可用回覆。

\subsection{量化指標}
\begin{itemize}[leftmargin=2em]
  \item 生成題有效率：\texttt{validatedCount/generatedCount}
  \item 錯誤碼分佈：\texttt{invalidReasonStats}
  \item 提示依賴層數：每題 \texttt{maxHintLayer}
  \item RAG 命中比例：\texttt{ragContextUsed} 與注入字數
  \item Socratic 重複度：回覆 n-gram overlap
\end{itemize}

\subsection{對照實驗建議}
建議以環境變數操控四類自變項：驗題開關、RAG 開關、驗題模型切換、Hint 推斷開關，搭配上述指標做探索性比較。若樣本有限，建議優先報告效應量與置信區間，而非僅 p 值。

\section{結論、限制與未來展望}
\subsection{研究結論}
本文完成一套可重現、可部署、可觀測的平面向量智能教學原型。從工程觀測可見，雙階段命題有助於降低部分生成錯誤，Hint 層級推斷可改善對話模板化，RAG 三層降級可提升失效情境下穩定性。整體而言，系統已具備作為教學輔助工具的原型價值。

\subsection{研究限制}
\begin{enumerate}[leftmargin=2em]
  \item 無帳號與跨裝置同步，長期學習軌跡不足。
  \item 題庫規模有限，部分 topic 多樣性不足。
  \item LLM 幻覺無法完全消除，只能透過工程護欄降低風險。
  \item 自動化測試尚不完整，回歸驗證仍依賴人工抽檢與 trace。
  \item 外部效度有限（單一場域與樣本）。
\end{enumerate}

\subsection{未來展望}
\begin{itemize}[leftmargin=2em]
  \item 導入帳號體系與持久化事件流，支援長期追蹤。
  \item 擴充題庫規模與多模態題型（圖形題、掃描題）。
  \item 進行多週期、較大樣本對照研究，驗證教育成效。
  \item 將框架擴展至其他高中數學單元，檢驗遷移性。
\end{itemize}

\section*{效度威脅}
\addcontentsline{toc}{section}{效度威脅}
\begin{itemize}[leftmargin=2em]
  \item \textbf{內部效度：}未採完整隨機分組與盲測，可能受學習動機差異影響。
  \item \textbf{外部效度：}樣本與課綱場域集中，跨地區推廣需再驗證。
  \item \textbf{建構效度：}部分指標（如問句重複度）屬代理變數，未必完整代表教學品質。
\end{itemize}

\section*{可重現性清單}
\addcontentsline{toc}{section}{可重現性清單}
\begin{enumerate}[leftmargin=2em]
  \item 提供 \texttt{.env.example} 與關鍵環境變數說明。
  \item 提供 Supabase SQL（pgvector extension 與 RPC）。
  \item 提供 Prompt 對照表與附錄原文。
  \item 提供 trace 啟用方式與樣本 JSONL。
  \item 固定 commit hash 與執行命令，確保版本可回放。
\end{enumerate}

\section*{網站截圖預留區（依 v2 風格）}
\addcontentsline{toc}{section}{網站截圖預留區（依 v2 風格）}
\begin{figure}[H]
  \centering
  \includegraphics[width=0.9\textwidth]{figures/page-diagnostic.png}
  \caption{網站頁面 A：診斷問卷（預留）}
\end{figure}

\begin{figure}[H]
  \centering
  \includegraphics[width=0.9\textwidth]{figures/page-practice.png}
  \caption{網站頁面 B：練習與導師對話（預留）}
\end{figure}

\begin{figure}[H]
  \centering
  \includegraphics[width=0.9\textwidth]{figures/page-mastery.png}
  \caption{網站頁面 C：掌握度儀表板（預留）}
\end{figure}

\section*{附錄一：核心 Prompt 原文（完整）}
\addcontentsline{toc}{section}{附錄一：核心 Prompt 原文（完整）}
\textbf{維護說明：}以下模板依目前程式碼整理，\texttt{\$\{...\}} 代表執行期字串模板代入。

\subsection*{A.1 診斷掌握分析 Prompt A（\texttt{lib/ai/prompt-a.ts}）}
\textbf{System}
\begin{verbatim}
你只輸出一段合法 JSON 物件，不要 markdown、不要程式碼围栏、不要任何 JSON 以外文字。
\end{verbatim}
\textbf{User}
\begin{verbatim}
你是一位澳門高中數學向量單元的資深教師助手。
學生剛完成了一份 8 題的初始診斷問卷，請根據作答結果分析其掌握程度。

===== 學生作答資料 =====
學生ID：${input.studentId}
作答時間：${input.totalTimeSeconds} 秒

題目結果（每題：題號 | 知識點 | 知識點標籤 | 難度 | 是否正確 | 耗時秒數；「知識點標籤」對應題庫與練習排序用的 topic 代碼，無則留空）：
${input.questionResultsJson}

===== 你的任務 =====
請生成一份 JSON 格式的掌握度分析報告，格式嚴格如下，不要輸出任何其他內容：

{
  "overall_level": "L1" | "L2" | "L3",
  "mastery_score": <0-100的整數，基於正確率與難度加權>,
  "strong_topics": [
    { "topic": "知識點名稱", "evidence": "一句話說明為什麼認為掌握好" }
  ],
  "weak_topics": [
    { "topic": "向量的模", "gap": "……", "topic_tag": "vec_magnitude" }
  ],
  "recommended_start_level": "L1" | "L2" | "L3",
  "student_summary": "用第二人稱（你）、溫和鼓勵的語氣，2-3句話向學生說明其目前狀況，不要說「根據分析」，要像老師跟學生說話。語言：繁體中文。",
  "teacher_note": "給系統後端的內部備注（英文），說明初始難度設置建議和需要重點關注的知識點，不展示給學生。"
}
\end{verbatim}

\subsection*{A.2 診斷填空題複核（\texttt{app/api/diagnostic/analyze/route.ts}）}
\textbf{System}
\begin{verbatim}
你是嚴謹的數學閱卷助理。只輸出 JSON 物件，不要 markdown，不要其他文字。
\end{verbatim}
\textbf{User}
\begin{verbatim}
請複核以下填空題作答。學生答案可能包含解題過程或多餘文字，請判斷其最終數學答案是否正確。

輸入（JSON）：
${JSON.stringify(fillPayload)}

請輸出：
{
  "reviews": [
    { "q": 題號, "correct": true/false }
  ]
}
\end{verbatim}

\subsection*{A.3 練習題 AI 排序（\texttt{prompt-practice-curation.ts}）}
\textbf{System}
\begin{verbatim}
你只輸出一段合法 JSON 物件，不要 markdown、不要程式碼围栏、不要 JSON 以外文字。
\end{verbatim}
\textbf{User}
\begin{verbatim}
你是澳門高中數學向量單元的題組編排助手。下面有一份「診斷掌握報告」以及目前題庫中可練習的題目清單（僅含 id、知識標籤、難度、題型、題幹摘要）。
===== 診斷報告（JSON）=====
${input.reportJson}
===== 總掌握摘要（JSON，可為 null）=====
${input.masteryJson ?? "null"}
===== 題目清單（JSON 陣列，每題一筆）=====
${input.catalogJson}
===== 題目 id 列表（與上表一致，供你核對）=====
${JSON.stringify(input.catalogIds)}
===== 任務 =====
請輸出 JSON，格式嚴格如下：
{
  "ordered_ids": ["題目id1", "題目id2", ...],
  "rationale": "用繁體中文簡短說明排序邏輯（給教學系統備查，不直接給學生看）"
}
\end{verbatim}

\subsection*{A.4 練習題 AI 生成（\texttt{prompt-practice-curation.ts}）}
\textbf{System}
\begin{verbatim}
你只輸出一段合法 JSON 物件，不要 markdown、不要程式碼围栏。JSON 內須含鍵 "questions"，其值為陣列。
\end{verbatim}
\textbf{User}
\begin{verbatim}
你是澳門高中數學向量單元的命題助手。下面有診斷報告、題庫中少量範例題，以及需要補強的 topic 標籤。
===== 診斷報告（JSON）=====
${input.reportJson}
===== 總掌握摘要（JSON，可為 null）=====
${input.masteryJson ?? "null"}
===== 題庫範例（JSON 陣列）=====
${input.sampleRowsJson}
===== RAG 檢索片段 =====
${input.ragContext?.trim() || "（本輪未取回 RAG 片段）"}
===== 建議優先命題的 topic 標籤 =====
${JSON.stringify(input.weakTopicTags.length ? input.weakTopicTags : ["（請綜合報告自行推斷）"])}
===== 任務 =====
請新命 ${input.generateCount} 題全新練習題，可混合選擇題與填空題，並輸出：
{
  "questions": [ { ...一題一物件... } ],
  "design_note": "繁體中文，說明命題取向（系統用）"
}
\end{verbatim}

\subsection*{A.5 練習題 AI 驗題（\texttt{prompt-practice-curation.ts}）}
\textbf{System}
\begin{verbatim}
你是嚴格的高中向量題品質審核器。你只輸出一段合法 JSON 物件，不要 markdown、不要程式碼围栏、不要額外文字。
\end{verbatim}
\textbf{User}
\begin{verbatim}
請審核以下 AI 生成題，找出邏輯或答案問題，並僅回傳 JSON。
===== 診斷報告（JSON）=====
${input.reportJson}
===== 總掌握摘要（JSON，可為 null）=====
${input.masteryJson ?? "null"}
===== 待審核題目（JSON）=====
${input.questionsJson}
===== 審核規則 =====
1) ANSWER_MISMATCH
2) OPTIONS_DUPLICATE
3) STEM_OPTION_CONTRADICTION
4) UNSOLVABLE_OR_NO_VALID_CHOICE
5) FORMAT_INVALID
\end{verbatim}

\subsection*{A.6 練習批次報告（\texttt{lib/ai/prompt-batch-report.ts}）}
\textbf{System}
\begin{verbatim}
你是澳門高中向量單元的學習分析助教。只輸出合法 JSON 物件，不要 markdown、不要程式碼區塊、不要其他文字。
\end{verbatim}
\textbf{User}
\begin{verbatim}
下面是學生在第 ${input.batchIndex + 1} 組（5題）練習的作答結果 JSON：
${JSON.stringify(input.results)}
===== 本組作答結構摘要（JSON）=====
${JSON.stringify(input.attemptsSummary ?? null)}
\end{verbatim}

\subsection*{A.7 總掌握報告（\texttt{lib/ai/prompts/mastery-report.ts}）}
\textbf{System}
\begin{verbatim}
你是高中向量單元的學習分析教練。你的任務是依據診斷與練習歷史，輸出一份可機械解析的掌握追蹤 JSON。
\end{verbatim}
\textbf{User}
\begin{verbatim}
請根據以下資料產生掌握追蹤 JSON。
[diagnostic]
${input.diagnosticJson}
[practice_events]
${input.practiceEventsJson}
[batch_reports]
${input.batchReportsJson}
[previous_report]
${input.previousReportJson}
\end{verbatim}

\subsection*{A.8 單題收束小結（\texttt{lib/ai/prompt-question-recap.ts}）}
\textbf{System}
\begin{verbatim}
你是高中數學向量單元導師。請用繁體中文寫一段精簡的本題小結，語氣溫和、第二人稱「你」。
\end{verbatim}
\textbf{User}
\begin{verbatim}
題幹要點（LaTeX 混排）：
${input.questionStemPlain}
知識向度：${input.topicLabel}
題型：${input.kind}
學生最終作答：${input.studentAnswerSummary}
標準參考：${input.standardAnswerSummary}
\end{verbatim}

\subsection*{A.9 Socratic 導師主對話（\texttt{lib/ai/prompt-b.ts}）}
\textbf{System}
\begin{verbatim}
你是澳門高中數學向量單元的 AI 教學助手「向量老師」。教學風格為 Socratic：不直接給出答案或完整解法，用提問引導。回應語言：繁體中文。
【數學式 LaTeX 規範】
- 行內公式：$...$
- 區塊公式：$$...$$
\end{verbatim}
\textbf{User}
\begin{verbatim}
===== 學生背景 =====
掌握程度等級：${input.studentLevel}
掌握良好的知識點：${JSON.stringify(input.strongTopics)}
需要加強的知識點：${JSON.stringify(input.weakTopics)}
當前總體掌握分：${input.masteryScore}
===== 當前題目資料 =====
題目編號：${input.questionId}
題目內容（LaTeX）：${input.questionLatex}
正確答案（僅供你參考，不得直接透露）：${input.correctAnswer}
===== 教學規則（不得違反）=====
1. 不直接說出正確答案或完整步驟
2. 每次回應以引導性問題結尾
3. 若 isCurrentAnswerCorrect=true，第一句先肯定
===== UI_SIGNAL =====
<!-- UI_SIGNAL -->
{
  "action": "next_question" | "show_hint_button" | "show_answer" | "encourage_retry" | "auto_reveal_hint",
  "update_mastery": true | false,
  "mastery_delta": 0,
  "topic_status": { "topic": "${input.questionTopic}", "result": "mastered" | "partial" | "needs_review" },
  "hint_index_used": -1
}
\end{verbatim}

\subsection*{A.10 提示改寫 Coach（\texttt{app/api/tutor/hint/route.ts}）}
\textbf{System}
\begin{verbatim}
你是澳門高中數學向量單元的導師。請把參考提示改寫成簡短引導（多為提問），不可連續複製原文超過 6 個字，不要直接給最終答案。
\end{verbatim}
\textbf{User}
\begin{verbatim}
參考提示（第 {hintIndex + 1} 層）：
{hintSourceText}

學生剛寫的思路（可為空）：
{studentThought 或 （無）}
\end{verbatim}

\section*{參考文獻}
\addcontentsline{toc}{section}{參考文獻}
\begin{enumerate}[leftmargin=2em]
  \item Nehring, J., Moyer-Packenham, P., \& North, M. (2023). Assessing the effectiveness of an AI tutoring system.
  \item Son, T. (2024). Intelligent Tutoring Systems in Mathematics Education.
  \item L\'etourneau, A. (2025). A systematic review of AI-driven ITS in K-12.
  \item Kestin, G., et al. (2025). AI tutoring outperforms active learning in RCT settings.
  \item Levonian, Z., et al. (2024). RAG trade-offs in math QA.
  \item Levonian, Z., et al. (2025). Designing safe and relevant generative chats for math learning.
  \item ECNU-ICALK. (2024). SocraticMath.
  \item Davis, F. D. (1989). Perceived usefulness and perceived ease of use.
\end{enumerate}

\end{document}
